\documentclass[10pt]{article}
\usepackage[margin=0.8in]{geometry}
\usepackage{amsmath,microtype}
\title{Guide: The Mechanical-Word Branch Is Complete}
\author{Collatz proof project}
\date{September 5, 2026}
\begin{document}
\maketitle
\section*{What Lean now proves}
For a positive natural seed, suppose the whole sequence of even and odd
orbit values is generated by
\[
 \lfloor(t+1)\alpha+\rho\rfloor-\lfloor t\alpha+\rho\rfloor.
\]
Then the seed is 1 or 2 and the slope is $\alpha=1/2$. This holds for every
real slope and intercept satisfying the itinerary hypothesis. It does not
mean every intercept works for those seeds.

If this mechanical pattern begins only after a finite orbit prefix, the
original seed still reaches 1. Zero is excluded from the hypothesis because
it has an all-zero pattern.

\section*{How the remaining rational case was closed}
A rational mechanical word is periodic. Equal infinite parity patterns
identify natural seeds exactly, so a parity period is an actual return
period for the Collatz orbit.

For a reduced slope $p/q$, the pattern can be coded by residues modulo $q$.
The residues visit every phase. Two extreme phases have period words of
the form $1u0$ and $0u1$: the middle is identical, but the endpoints swap.
The previous Lean theorem says both words can be natural returns only when
$u$ is empty, with seeds 1 and 2.

Every real intercept has an exact residue phase. This covers negative
intercepts and values on the floor thresholds, without numerical
approximation. Reducing a fraction by its greatest common divisor handles
nonprimitive presentations as well.

\section*{Why this covers arbitrary real slopes}
The earlier complexity theorem forces any mechanical natural orbit to be
bounded. A bounded orbit repeats a state. Counting odd steps over repeated
periods then forces the slope to be rational, and the rational result gives
a visit to 1. That visit forces slope $1/2$; its period of two applies back
at the original seed, which must itself be 1 or 2.

\section*{What remains open}
This completes the mechanical itinerary class, extending the earlier
individual slopes, infinite slope families, and irrational-only exclusion.
It is a formal integration of classical ideas, not a claim of mathematical
novelty or a proof of Collatz.

There is no proof that arbitrary positive Collatz orbits eventually have
mechanical parity patterns. Assuming that would hide the main difficulty.
Future work should address general counterexample restrictions or descent,
instead of adding more mechanical slope exclusions.

\section*{Verification and files}
The Lean build passes; 107 selected declarations have only standard
foundational axiom dependencies. Four exact finite regression controls pass.
The universal statement comes from Lean, not those finite controls.
Run \texttt{lake build Collatz.RationalMechanical} from \texttt{lean/}.
The technical paper and editable LaTeX sources are in the paper directory.
\end{document}
